% amortized_inference.tex — Amortized Neural Inference for Seismic Inversion
% Compile with: lualatex amortized_inference.tex
\documentclass[11pt,a4paper]{article}

\usepackage{fontspec}
\setmainfont{Latin Modern Roman}
\setmonofont{Latin Modern Mono}

\usepackage{amsmath,amssymb,amsthm}
\usepackage{mathtools}
\usepackage{bm}
\usepackage{booktabs}
\usepackage[svgnames]{xcolor}
\usepackage{geometry}
\geometry{margin=2.5cm}

\usepackage{hyperref}
\hypersetup{
  colorlinks=true,
  linkcolor=DarkBlue,
  citecolor=DarkBlue,
  urlcolor=DarkBlue,
}

\usepackage{enumitem}
\usepackage{fancyhdr}
\pagestyle{fancy}
\fancyhf{}
\fancyhead[L]{\small Amortized Neural Inference for Seismic Inversion}
\fancyhead[R]{\small\thepage}
\renewcommand{\headrulewidth}{0.4pt}

% Notation shortcuts
\newcommand{\dd}{\mathrm{d}}
\newcommand{\ppd}[2]{\frac{\partial #1}{\partial #2}}
\newcommand{\KL}{\mathrm{KL}}
\newcommand{\ELBO}{\mathcal{L}_{\mathrm{ELBO}}}
\newcommand{\Robs}{R_{\mathrm{obs}}}
\newcommand{\xobs}{\mathbf{x}_{\mathrm{obs}}}

\theoremstyle{plain}
\newtheorem{proposition}{Proposition}
\theoremstyle{remark}
\newtheorem{remark}{Remark}

\title{%
  Amortized Neural Inference for Seismic Inversion\\[6pt]
  \large Transformer Architecture, Training Procedure,\\
  and Hybrid Newton-Net Pipeline%
}
\author{}
\date{}

\begin{document}
\maketitle
\thispagestyle{fancy}

\begin{abstract}
We describe an amortized variational inference approach to the seismic
inverse problem in stratified elastic media.  A Transformer-based
inference network, \texttt{SeismicInferenceNet}, learns a mapping from
plane-wave reflectivity $R(\omega, p)$ to a diagonal Gaussian
approximate posterior over log-space earth model parameters.  Once
trained, the network provides instant posterior estimates for any new
observation, eliminating the iterative cost of classical
Levenberg--Marquardt inversion.

The network prediction serves as a warm start for a subsequent
Newton--LM refinement stage (the ``hybrid Newton-Net pipeline''),
combining the amortized posterior with the quadratic convergence of the
exact-Hessian Newton method.  We analyse the model-specificity of the
trained network, identify what changes in the forward model require
retraining, and derive the break-even formula for when the hybrid
approach is cost-effective relative to random-start Newton inversion.
\end{abstract}

\tableofcontents

%=============================================================================
\section{Amortized Variational Inference}\label{sec:amortized-vi}
%=============================================================================

\subsection{The inverse problem}

The seismic inverse problem seeks to recover the earth model parameter
vector
\begin{equation}\label{eq:params}
  \mathbf{m}
  = \bigl(\alpha_1,\dotsc,\alpha_N,\;
          \beta_1,\dotsc,\beta_N,\;
          \rho_1,\dotsc,\rho_N,\;
          h_1,\dotsc,h_{N-1}\bigr)
  \in \mathbb{R}^{4N-1}
\end{equation}
from observed plane-wave reflectivity $\Robs(\omega, p)$, where
$\alpha_j, \beta_j, \rho_j, h_j$ are the P-wave velocity, S-wave
velocity, density, and thickness of layer~$j$.

The forward operator $\mathcal{F}: \mathbf{m} \mapsto R(\omega, p)$
is computed by either the Kennett reflectivity recursion or the
Global Matrix Method (see companion document).  The problem is
ill-posed: band-limited data, velocity--thickness trade-offs, and
noise mean that many models fit the observations equally well.
A single ``best-fit'' model is therefore incomplete ---
\emph{uncertainty quantification} is essential.

\subsection{Classical approach: Newton--LM iteration}

The deterministic approach minimises the L2 misfit
\begin{equation}\label{eq:misfit}
  \chi(\mathbf{m})
  = \sum_{s,i}
    \bigl|R(\omega_i, p_s; \mathbf{m})
          - \Robs(\omega_i, p_s)\bigr|^2
\end{equation}
via Levenberg--Marquardt iteration with exact Hessian.  This has
quadratic convergence near the solution but requires:
\begin{itemize}[nosep]
  \item a good starting model (convergence is local),
  \item Jacobian and Hessian computation at every iteration,
  \item a separate optimisation for each observation.
\end{itemize}

\subsection{Amortized approach}

Instead of solving an optimisation problem per observation, we train a
neural network $q_\phi(\mathbf{m} \mid \mathbf{x})$ that maps any
observation $\mathbf{x} = R(\omega, p)$ directly to an approximate
posterior distribution over~$\mathbf{m}$.  The key distinction is that
the network learns a \textbf{mapping from data to posterior}, not a
single point solution --- the upfront training cost is incurred once,
and inference for any new observation is a single forward pass.

Formally, we seek to minimise the expected KL divergence between the
approximate posterior and the true posterior:
\begin{equation}\label{eq:amortized-kl}
  \min_\phi\;
  \mathbb{E}_{p(\mathbf{x})}
  \bigl[\KL\!\bigl(q_\phi(\mathbf{m}\mid\mathbf{x})
    \,\|\, p(\mathbf{m}\mid\mathbf{x})\bigr)\bigr],
\end{equation}
where $p(\mathbf{x})$ is the marginal data distribution induced by the
prior $p(\mathbf{m})$ and forward model $\mathcal{F}$.

\subsection{Diagonal Gaussian NLL as ELBO}

We parameterise the approximate posterior as a diagonal Gaussian in
log-parameter space:
\begin{equation}\label{eq:approx-posterior}
  q_\phi(\mathbf{z}\mid\mathbf{x})
  = \prod_{k=1}^{P} \mathcal{N}\bigl(z_k;\; \mu_k(\mathbf{x}),\;
    \sigma_k^2(\mathbf{x})\bigr),
  \qquad \mathbf{z} = \log\mathbf{m},
\end{equation}
where $\mu_k$ and $\sigma_k$ are output by the network.  The training
loss is the diagonal Gaussian negative log-likelihood:
\begin{equation}\label{eq:nll-loss}
  \mathcal{L}(\phi)
  = \mathbb{E}_{(\mathbf{x},\mathbf{z})\sim p(\mathbf{x},\mathbf{z})}
    \left[
      \sum_{k=1}^{P}
      \left(
        \frac{(z_k - \mu_k)^2}{2\sigma_k^2} + \log\sigma_k
      \right)
    \right],
\end{equation}
which corresponds to the negative ELBO under amortized variational
inference with a simulation-based prior.  Training samples
$(\mathbf{x}, \mathbf{z})$ are generated by drawing
$\mathbf{m} \sim p(\mathbf{m})$, computing
$\mathbf{x} = \mathcal{F}(\mathbf{m})$, and setting
$\mathbf{z} = \log\mathbf{m}$.

%=============================================================================
\section{SeismicInferenceNet Architecture}\label{sec:architecture}
%=============================================================================

\subsection{Input representation}

The observed reflectivity $R(\omega, p)$ is complex-valued with shape
$(n_p \times n_\omega)$.  We stack the real and imaginary parts along
the feature dimension to produce a real-valued input tensor:
\begin{equation}\label{eq:input}
  \mathbf{X} \in \mathbb{R}^{n_p \times 2n_\omega},
  \qquad
  \mathbf{X}_{j,:} = \bigl[\operatorname{Re} R(\cdot, p_j),\;
                            \operatorname{Im} R(\cdot, p_j)\bigr].
\end{equation}
Each slowness value $p_j$ is treated as a token with the spectrum as
its feature vector.  Per-sample normalisation (dividing by
$\max|\mathbf{X}|$) ensures scale invariance.

\subsection{Input projection}

A linear layer projects each token from $2n_\omega$ features to
$d_{\text{model}}$ dimensions:
\begin{equation}\label{eq:input-proj}
  \mathbf{h}_j^{(0)}
  = \mathbf{W}_{\text{proj}}\,\mathbf{X}_{j,:} + \mathbf{b}_{\text{proj}},
  \qquad
  \mathbf{W}_{\text{proj}} \in \mathbb{R}^{d_{\text{model}} \times 2n_\omega}.
\end{equation}

\subsection{Continuous positional encoding}

Unlike standard Transformer architectures that use integer position
indices, \texttt{SeismicInferenceNet} encodes the \emph{physical}
slowness value $p_j$ via sinusoidal features:
\begin{equation}\label{eq:pos-enc}
  \text{PE}(p_j)
  = \bigl[\sin(f_1 p_j),\; \cos(f_1 p_j),\;
          \sin(f_2 p_j),\; \cos(f_2 p_j),\;
          \dotsc\bigr]
  \in \mathbb{R}^{d_{\text{model}}},
\end{equation}
where $\{f_k\}$ are geometrically spaced frequencies from $1$ to
$f_{\max}$.  The positional encoding is added to the projected tokens:
$\mathbf{h}_j^{(0)} \leftarrow \mathbf{h}_j^{(0)} + \text{PE}(p_j)$.

\begin{remark}
Continuous positional encoding enables \textbf{generalisation to unseen
slowness grids} without retraining.  A network trained on 12 uniformly
spaced slownesses can be applied to data acquired at arbitrary ray
parameters, since the encoding is a smooth function of~$p$.
\end{remark}

\subsection{Transformer encoder}

The projected tokens are processed by a stack of $L$ pre-norm
Transformer encoder layers.  Each layer applies:
\begin{enumerate}[nosep]
  \item \textbf{Layer normalisation} (pre-norm for stable training),
  \item \textbf{Multi-head self-attention} across slowness tokens ---
    this learns cross-slowness (AVO) correlations,
  \item \textbf{Feed-forward network} with GELU activation.
\end{enumerate}
Self-attention over the slowness dimension allows the network to learn
how reflectivity varies with ray parameter --- precisely the
amplitude-versus-offset (AVO) information that constrains the elastic
parameters.

\subsection{Output heads}

After the Transformer encoder, the token representations are
mean-pooled over the slowness dimension:
\begin{equation}\label{eq:mean-pool}
  \bar{\mathbf{h}} = \frac{1}{n_p}\sum_{j=1}^{n_p} \mathbf{h}_j^{(L)}
  \in \mathbb{R}^{d_{\text{model}}}.
\end{equation}
Two separate linear heads produce the posterior parameters:
\begin{align}
  \bm{\mu} &= \mathbf{W}_\mu\,\bar{\mathbf{h}} + \mathbf{b}_\mu
    \in \mathbb{R}^P, \label{eq:mu-head}\\
  \log\bm{\sigma} &= \operatorname{clamp}\!\bigl(
    \mathbf{W}_\sigma\,\bar{\mathbf{h}} + \mathbf{b}_\sigma,\;
    -5,\; 2\bigr)
    \in \mathbb{R}^P, \label{eq:sigma-head}
\end{align}
where $P = 4N - 1$ is the number of model parameters and the clamping
prevents numerical instability.  The bias of the $\log\sigma$ head is
initialised to $\log(0.1) \approx -2.3$ for stable early training.

\subsection{Architecture hyperparameters}

\begin{center}
\renewcommand{\arraystretch}{1.3}
\begin{tabular}{llr}
  \toprule
  \textbf{Hyperparameter} & \textbf{Symbol} & \textbf{Value} \\
  \midrule
  Embedding dimension     & $d_{\text{model}}$ & 64 \\
  Attention heads         & $n_{\text{heads}}$ & 4 \\
  Encoder layers          & $L$                & 3 \\
  Feed-forward dimension  & $d_{\text{ff}}$    & 128 \\
  Dropout rate            & ---                & 0.1 \\
  Positional encoding     & ---                & continuous sinusoidal \\
  Activation              & ---                & GELU \\
  Normalisation           & ---                & pre-norm \\
  \bottomrule
\end{tabular}
\end{center}

%=============================================================================
\section{Training Procedure}\label{sec:training}
%=============================================================================

\subsection{Data generation}

Training data is generated by sampling earth models from a prior
distribution centred on a reference ocean-crust model.  For each
sample:
\begin{enumerate}[nosep]
  \item Draw multiplicative perturbation factors from uniform
    distributions: velocities $\mathcal{U}(1{-}\delta_v,\; 1{+}\delta_v)$,
    densities $\mathcal{U}(1{-}\delta_\rho,\; 1{+}\delta_\rho)$,
    thicknesses $\mathcal{U}(1{-}\delta_h,\; 1{+}\delta_h)$.
  \item Apply perturbations to the reference parameter vector.
  \item Compute reflectivity $R(\omega, p)$ using the GMM
    Riccati-sweep solver for all slownesses and frequencies.
  \item Store $(\mathbf{X}, \mathbf{z}) = (\text{reflectivity},\;
    \log\mathbf{m})$ as a training pair.
\end{enumerate}

\subsection{Optimiser and schedule}

The network is trained with the AdamW optimiser and a cosine annealing
learning rate schedule.  Training terminates either at the maximum
epoch count or when validation loss has not improved for
$P_{\text{patience}}$ consecutive epochs (early stopping).  The best
checkpoint (lowest validation NLL) is saved and restored at the end.

\subsection{Training hyperparameters}

\begin{center}
\renewcommand{\arraystretch}{1.3}
\begin{tabular}{llr}
  \toprule
  \textbf{Hyperparameter} & \textbf{Symbol} & \textbf{Value} \\
  \midrule
  Training samples    & $N_{\text{train}}$    & 50{,}000 \\
  Validation samples  & $N_{\text{val}}$      & 5{,}000 \\
  Test samples        & $N_{\text{test}}$     & 5{,}000 \\
  Maximum epochs      & $E_{\max}$            & 200 \\
  Batch size          & $B$                   & 256 \\
  Learning rate       & $\eta$                & $10^{-3}$ \\
  Weight decay        & $\lambda_w$           & $10^{-4}$ \\
  Early stopping patience & $P_{\text{patience}}$ & 20 \\
  Gradient clipping   & $\|\nabla\|_{\max}$   & 1.0 \\
  Velocity perturbation  & $\delta_v$         & 20\% \\
  Density perturbation   & $\delta_\rho$      & 15\% \\
  Thickness perturbation & $\delta_h$         & 25\% \\
  Slowness grid       & $n_p$                 & 12 \\
  Frequency bins      & $n_\omega$            & 64 \\
  \bottomrule
\end{tabular}
\end{center}

\subsection{Evaluation metrics}

The trained network is evaluated on a held-out test set using three
metrics:
\begin{itemize}[nosep]
  \item \textbf{Test NLL}: the diagonal Gaussian negative
    log-likelihood~\eqref{eq:nll-loss} on test samples.
  \item \textbf{Mean absolute error}: $\mathbb{E}[|\mu_k - z_k|]$
    averaged over all parameters and test samples.
  \item \textbf{Calibration}: fraction of true parameters falling
    within the predicted $2\sigma$ interval (should be ${\approx}95\%$
    for a well-calibrated Gaussian).
\end{itemize}

\noindent Results on the 4-layer ocean-crust model:
\begin{center}
\renewcommand{\arraystretch}{1.3}
\begin{tabular}{lr}
  \toprule
  \textbf{Metric} & \textbf{Value} \\
  \midrule
  Test NLL (nats/param) & $-62.69$ \\
  Mean absolute error (log-space) & 0.0083 \\
  Calibration ($2\sigma$) & 99.5\% \\
  Best epoch & 194 / 200 \\
  \bottomrule
\end{tabular}
\end{center}

\noindent The high calibration ($99.5\%$ vs.\ the nominal $95.4\%$)
indicates slightly conservative uncertainty estimates, which is
desirable for warm-starting Newton iterations.

%=============================================================================
\section{Model-Specificity and Retraining}\label{sec:specificity}
%=============================================================================

The trained network is \textbf{specialised to one model class}.  A
``model class'' is defined by the combination of:
\begin{enumerate}[nosep]
  \item \textbf{Layer count} $N$: determines the output dimension
    $P = 4N - 1$.
  \item \textbf{Parameter ordering}: the mapping from packed vector
    indices to physical quantities.
  \item \textbf{Prior distribution}: the reference model and
    perturbation ranges used for training data generation.
  \item \textbf{Frequency grid}: the number of frequency bins $n_\omega$
    determines the input projection width.
\end{enumerate}

\subsection{Changes that require retraining}

\begin{center}
\renewcommand{\arraystretch}{1.3}
\begin{tabular}{lp{8cm}}
  \toprule
  \textbf{Change} & \textbf{Reason} \\
  \midrule
  Different $N$ (layer count)
    & Output heads have dimension $P = 4N - 1$; the network architecture
      itself changes. \\
  Different geology (prior shift)
    & Training data no longer represents the target distribution;
      the network's learned mapping is biased. \\
  Different $n_\omega$ (frequency bins)
    & Input projection $\mathbf{W}_{\text{proj}} \in
      \mathbb{R}^{d_{\text{model}} \times 2n_\omega}$ has a different
      width. \\
  \bottomrule
\end{tabular}
\end{center}

\subsection{Changes that do NOT require retraining}

\begin{center}
\renewcommand{\arraystretch}{1.3}
\begin{tabular}{lp{8cm}}
  \toprule
  \textbf{Change} & \textbf{Reason} \\
  \midrule
  Different slowness grid
    & Continuous positional encoding~\eqref{eq:pos-enc} generalises to
      arbitrary $p$ values; the network processes each slowness token
      independently before attention. \\
  Different data observations
    & This is the entire point of amortization --- the network maps
      \emph{any} observation to a posterior without re-optimisation. \\
  \bottomrule
\end{tabular}
\end{center}

\begin{remark}
The model-specificity of amortized inference is the fundamental
trade-off: the network provides instant posteriors but only for the
model class it was trained on.  This contrasts with classical Newton
inversion, which adapts to any parameterisation but requires iterative
optimisation for every observation.
\end{remark}

%=============================================================================
\section{Hybrid Newton-Net Pipeline}\label{sec:hybrid}
%=============================================================================

The hybrid pipeline combines the strengths of both approaches:
the network provides a good starting point with calibrated uncertainty,
and Newton--LM refines it to high precision.

\subsection{Two-stage algorithm}

\begin{enumerate}
  \item \textbf{Stage~1: Network prediction.}  Feed the observed
    reflectivity $\Robs(\omega, p)$ through the trained network to
    obtain $(\bm{\mu}, \bm{\sigma})$ --- a warm start and calibrated
    uncertainty in log-parameter space.  Cost: a single forward pass
    (${\sim}\text{ms}$).

  \item \textbf{Stage~2: Newton--LM refinement.}  Starting from
    $\mathbf{z}^{(0)} = \bm{\mu}$, iterate the Levenberg--Marquardt
    update
    \begin{equation}\label{eq:lm-update}
      (\mathbf{H} + \lambda\mathbf{I})\,\delta\mathbf{z}
      = -\mathbf{g}
    \end{equation}
    with exact Hessian $\mathbf{H}$ and adaptive damping until
    convergence.  The network's warm start places the initial guess
    close to the solution.
\end{enumerate}

\subsection{Why this works}

Newton's method has \textbf{quadratic convergence near the solution}:
$\|\mathbf{z}_{k+1} - \mathbf{z}^*\| \leq c\|\mathbf{z}_k -
\mathbf{z}^*\|^2$.  But this convergence is only guaranteed within
a neighbourhood of the true solution --- far from it, Newton may
diverge or converge to a local minimum.

The network's prediction has low relative error (${\sim}0.2\%$ on the
reference model), placing the initial guess firmly within the quadratic
convergence basin.  The Newton refinement then achieves machine-precision
accuracy in few iterations.

\subsection{Convergence comparison}

Results on the 4-layer ocean-crust model ($P = 15$ parameters):

\begin{center}
\renewcommand{\arraystretch}{1.3}
\begin{tabular}{lrr}
  \toprule
  & \textbf{Hybrid} & \textbf{Random start} \\
  \midrule
  Iterations          & 10        & 28 \\
  Converged           & yes       & yes \\
  Final misfit $\chi$ & $2.61\times 10^{-7}$ & $1.42\times 10^{-22}$ \\
  Final parameter error & $2.82\times 10^{-5}$ & $1.06\times 10^{-12}$ \\
  Wall time (s)       & 28.7      & 80.9 \\
  \bottomrule
\end{tabular}
\end{center}

\noindent The hybrid approach uses $2.8\times$ fewer iterations and
$2.8\times$ less wall time.  The random-start Newton achieves higher
final precision because it runs more iterations, but both methods
converge to the true model.  If higher precision is needed, the hybrid
pipeline can simply run additional Newton iterations from its already
excellent starting point.

\subsection{Network prediction quality}

Before any Newton refinement, the network's prediction achieves:

\begin{center}
\renewcommand{\arraystretch}{1.3}
\begin{tabular}{lr}
  \toprule
  \textbf{Metric} & \textbf{Value} \\
  \midrule
  Relative parameter error & $2.20 \times 10^{-3}$ \\
  Max single-parameter error & $< 1\%$ \\
  \bottomrule
\end{tabular}
\end{center}

\noindent This sub-percent accuracy from a single forward pass ---
without any iterative optimisation --- demonstrates the power of
amortized inference.

%=============================================================================
\section{Amortization Cost Analysis}\label{sec:cost}
%=============================================================================

\subsection{Upfront cost}

The amortized approach requires a one-time investment:

\begin{center}
\renewcommand{\arraystretch}{1.3}
\begin{tabular}{lr}
  \toprule
  \textbf{Stage} & \textbf{Wall time (s)} \\
  \midrule
  Data generation (50{,}000 samples $\times$ 12 slownesses) & 4{,}747 \\
  Network training (200 epochs, early stop at 194) & 3{,}322 \\
  \midrule
  \textbf{Total upfront cost} $C_{\text{up}}$ & \textbf{8{,}069} \\
  \bottomrule
\end{tabular}
\end{center}

\subsection{Per-inversion savings}

For each subsequent inversion with the same model class:
\begin{align}
  T_{\text{random}} &= 80.9\;\text{s} && \text{(random-start Newton)},
  \label{eq:t-random}\\
  T_{\text{hybrid}} &= 28.7\;\text{s} && \text{(network + Newton)},
  \label{eq:t-hybrid}\\
  \Delta T &= T_{\text{random}} - T_{\text{hybrid}}
            = 52.2\;\text{s/inversion}. \label{eq:delta-t}
\end{align}

\subsection{Break-even formula}

The hybrid approach becomes cost-effective after $N_{\text{break}}$
inversions:
\begin{equation}\label{eq:breakeven}
  \boxed{
  N_{\text{break}}
  = \left\lceil \frac{C_{\text{up}}}{\Delta T} \right\rceil
  = \left\lceil \frac{C_{\text{up}}}{T_{\text{random}} - T_{\text{hybrid}}}
    \right\rceil.
  }
\end{equation}

\noindent For the 4-layer ocean-crust model:
\begin{equation}\label{eq:breakeven-value}
  N_{\text{break}}
  = \left\lceil \frac{8{,}069}{52.2} \right\rceil
  = 155\;\text{inversions}.
\end{equation}

\subsection{When amortization pays off}

\begin{center}
\renewcommand{\arraystretch}{1.3}
\begin{tabular}{p{6cm}p{6cm}}
  \toprule
  \textbf{Favourable scenarios} & \textbf{Unfavourable scenarios} \\
  \midrule
  Many inversions with the same model class: marine surveys with
  hundreds of CMP gathers, time-lapse monitoring with repeated
  acquisitions.
  &
  One-off inversions with unique parameterisations: exploratory
  studies, models with varying layer counts. \\[6pt]
  Need for calibrated uncertainty: the network provides a free
  approximate posterior alongside the point estimate.
  &
  Rapidly changing model structure: if the geology requires
  different $N$ at every location, the retraining cost dominates. \\[6pt]
  Multi-modal posteriors: random-start Newton may converge to a
  local minimum, while the network is trained on the full prior
  and can navigate the posterior landscape.
  &
  Extreme precision required with few inversions: random-start
  Newton with sufficient iterations achieves higher final
  precision. \\
  \bottomrule
\end{tabular}
\end{center}

\begin{remark}
The break-even analysis considers only wall time.  Additional benefits
of amortization --- calibrated uncertainty, robustness to local minima,
and instant posterior estimates for real-time applications --- are not
captured by the simple cost formula~\eqref{eq:breakeven}.
\end{remark}

%=============================================================================
\section*{References}
%=============================================================================

\begin{enumerate}[label={[\arabic*]},nosep]
  \item D.P.\ Kingma and M.\ Welling, ``Auto-encoding variational Bayes,''
    \emph{arXiv:1312.6114}, 2013.
  \item K.\ Cranmer, J.\ Brehmer and G.\ Louppe, ``The frontier of
    simulation-based inference,'' \emph{Proc.\ Natl.\ Acad.\ Sci.},
    vol.~117, no.~48, pp.~30055--30062, 2020.
  \item A.\ Siahkoohi, G.\ Rizzuti, M.\ Louboutin, P.A.\ Witte and
    F.J.\ Herrmann, ``Deep-learning based Bayesian inference for seismic
    imaging,'' \emph{EAGE}, arXiv:2001.04567, 2020.
  \item A.\ Siahkoohi, G.\ Rizzuti and F.J.\ Herrmann, ``Amortized
    variational Bayesian inference with conditional normalizing flows,''
    \emph{EAGE}, arXiv:2203.15881, 2022.
  \item A.\ Spurio Mancini et al., ``Neural posterior estimation in
    seismology,'' \emph{Geophys.\ J.\ Int.}, vol.~241, no.~3, 2025.
  \item A.\ Vaswani et al., ``Attention is all you need,''
    \emph{Advances in Neural Information Processing Systems}, vol.~30,
    2017.
  \item D.W.\ Marquardt, ``An algorithm for least-squares estimation of
    nonlinear parameters,'' \emph{SIAM J.\ Appl.\ Math.}, vol.~11,
    pp.~431--441, 1963.
  \item B.L.N.\ Kennett, \emph{Seismic Wave Propagation in Stratified
    Media}, Cambridge University Press, 1983.
  \item I.\ Loshchilov and F.\ Hutter, ``Decoupled weight decay
    regularization,'' \emph{arXiv:1711.05101}, 2017.
\end{enumerate}

\end{document}
